\subsection*{Algebra}
\unskip

%=========================================================
% TRIGONOMETRY
%=========================================================

\begin{mytikzbox}{Trigonometry}{
\textbf{Angle addition:}
\begin{gather*}
\sin(a\pm b)
=
\sin a\cos b\pm\cos a\sin b,
\\
\cos(a\pm b)
=
\cos a\cos b\mp\sin a\sin b.
\end{gather*}

\textbf{Product-to-sum identities:}
\begin{gather*}
\cos a\cos b
=
\frac{1}{2}
\left[
\cos(a+b)+\cos(a-b)
\right],
\\
\sin a\sin b
=
\frac{1}{2}
\left[
\cos(a-b)-\cos(a+b)
\right],
\\
\cos a\sin b
=
\frac{1}{2}
\left[
\sin(a+b)-\sin(a-b)
\right].
\end{gather*}

\textbf{Complex representation:}
\begin{gather*}
\sin z
=
\frac{e^{iz}-e^{-iz}}{2i},
\qquad
\cos z
=
\frac{e^{iz}+e^{-iz}}{2},
\\
\cosh x=\cos(ix),
\qquad
\sinh x=-i\sin(ix).
\end{gather*}

\textbf{Double-angle identities:}
\[
\cos(2x)=\cos^2x-\sin^2x,
\qquad
\sin(2x)=2\sin x\cos x.
\]

\textbf{Half-angle identities:}
\begin{gather*}
\sin\left(\frac{\theta}{2}\right)
=
\pm\sqrt{\frac{1-\cos\theta}{2}},
\qquad
\cos\left(\frac{\theta}{2}\right)
=
\pm\sqrt{\frac{1+\cos\theta}{2}},
\\
\tan\left(\frac{\theta}{2}\right)
=
\pm\sqrt{\frac{1-\cos\theta}{1+\cos\theta}}
=
\frac{\sin\theta}{1+\cos\theta}
=
\frac{1-\cos\theta}{\sin\theta}.
\end{gather*}
\vspace{-10pt}
}{0.46}
\end{mytikzbox}

%=========================================================

\begin{mytikzbox}{Powers of trigonometric functions}{
\vspace{-7pt}
\begin{gather*}
\cos^2x
=
\frac{1+\cos(2x)}{2},
\qquad
\sin^2x
=
\frac{1-\cos(2x)}{2},
\\
\cos^3x
=
\frac{3}{4}\cos x+\frac{1}{4}\cos(3x),
\\
\sin^3x
=
\frac{3}{4}\sin x-\frac{1}{4}\sin(3x).
\end{gather*}
\vspace{-12pt}
}{0.46}
\end{mytikzbox}

%=========================================================
% FUNCTIONAL DERIVATIVES
%=========================================================

\begin{mytikzbox}{Functional derivatives}{
For
\[
J[f]
=
\int_a^b
\mathcal L\!\left(x,f(x),f'(x)\right)\,dx,
\]
the functional derivative is
\[
\frac{\delta J}{\delta f(x)}
=
\frac{\partial\mathcal L}{\partial f}
-
\frac{d}{dx}
\left(
\frac{\partial\mathcal L}{\partial f'}
\right).
\]

More generally, for a field $\phi(\mathbf x)$,
\[
J[\phi]
=
\int d^d x\,
\mathcal L(\phi,\partial_i\phi),
\]
\[
\frac{\delta J}{\delta\phi(\mathbf x)}
=
\frac{\partial\mathcal L}{\partial\phi}
-
\partial_i
\left(
\frac{\partial\mathcal L}
{\partial(\partial_i\phi)}
\right).
\]

Useful identity:
\[
\frac{\delta\phi(\mathbf x)}
{\delta\phi(\mathbf y)}
=
\delta^{(d)}(\mathbf x-\mathbf y).
\]
}{0.46}
\end{mytikzbox}

%=========================================================
% SPHERICAL COORDINATES
%=========================================================

\begin{mytikzbox}{Spherical coordinates}{
\textbf{Coordinate transformation:}
\begin{gather*}
x=r\sin\theta\cos\phi,
\qquad
y=r\sin\theta\sin\phi,
\\
z=r\cos\theta.
\end{gather*}

\textbf{Ranges:}
\[
r\geq0,
\qquad
0\leq\theta\leq\pi,
\qquad
0\leq\phi<2\pi.
\]

\textbf{Volume element:}
\[
d^3\mathbf r
=
r^2\sin\theta\,dr\,d\theta\,d\phi.
\]

\textbf{Radially symmetric functions:}
\[
\nabla f(r)
=
f'(r)\hat{\mathbf r},
\]
\[
\nabla^2 f(r)
=
\frac{1}{r^2}
\frac{d}{dr}
\left(
r^2\frac{df}{dr}
\right).
\]

For $\mathbf A=A_r(r)\hat{\mathbf r}$,
\[
\nabla\cdot\mathbf A
=
\frac{1}{r^2}
\frac{d}{dr}\left(r^2A_r\right).
\]
}{0.46}
\end{mytikzbox}

%=========================================================
% INDEFINITE INTEGRALS
%=========================================================

\begin{mytikzbox}{Indefinite integrals}{
\vspace{-6pt}
\begin{gather*}
\int xe^{ax}\,dx
=
\frac{e^{ax}}{a}
\left(
x-\frac{1}{a}
\right),
\\
\int x^2e^{ax}\,dx
=
\frac{e^{ax}}{a}
\left(
x^2-\frac{2x}{a}+\frac{2}{a^2}
\right),
\\
\int x\cos(ax)\,dx
=
\frac{x\sin(ax)}{a}
+
\frac{\cos(ax)}{a^2},
\\
\int x\sin(ax)\,dx
=
-\frac{x\cos(ax)}{a}
+
\frac{\sin(ax)}{a^2},
\\
\int x^2\cos(ax)\,dx
=
\frac{x^2\sin(ax)}{a}
+
\frac{2x\cos(ax)}{a^2}
-
\frac{2\sin(ax)}{a^3},
\\
\int x^2\sin(ax)\,dx
=
-\frac{x^2\cos(ax)}{a}
+
\frac{2x\sin(ax)}{a^2}
+
\frac{2\cos(ax)}{a^3}.
\end{gather*}

\begin{gather*}
\int\sin(ax)\cos(bx)\,dx
=
-\frac{\cos[(a+b)x]}{2(a+b)}
-\frac{\cos[(a-b)x]}{2(a-b)},
\\
\int\cos(ax)\cos(bx)\,dx
=
\frac{\sin[(a-b)x]}{2(a-b)}
+
\frac{\sin[(a+b)x]}{2(a+b)},
\\
\int\sin(ax)\sin(bx)\,dx
=
\frac{\sin[(a-b)x]}{2(a-b)}
-
\frac{\sin[(a+b)x]}{2(a+b)}.
\end{gather*}

\begin{gather*}
\int\sin^n(ax)\cos(ax)\,dx
=
\frac{\sin^{n+1}(ax)}{(n+1)a},
\\
\int\cos^n(ax)\sin(ax)\,dx
=
-\frac{\cos^{n+1}(ax)}{(n+1)a},
\\
\int e^{ax}\sin(bx)\,dx
=
\frac{e^{ax}}{a^2+b^2}
\left[
a\sin(bx)-b\cos(bx)
\right],
\\
\int e^{ax}\cos(bx)\,dx
=
\frac{e^{ax}}{a^2+b^2}
\left[
a\cos(bx)+b\sin(bx)
\right].
\end{gather*}
\vspace{-13pt}
}{0.46}
\end{mytikzbox}

%=========================================================
% DEFINITE INTEGRALS
%=========================================================

\begin{mytikzbox}{Definite integrals}{
\vspace{-7pt}
For $\Re(\alpha)>0$,
\begin{gather*}
\int_{-\infty}^{\infty}
e^{-\alpha x^2}\,dx
=
\sqrt{\frac{\pi}{\alpha}},
\\
\int_{-\infty}^{\infty}
e^{-\alpha x^2+i\beta x}\,dx
=
\sqrt{\frac{\pi}{\alpha}}
e^{-\beta^2/(4\alpha)}.
\end{gather*}

For $a>0$,
\begin{gather*}
\int_0^\infty
e^{-ax}\cos(bx)\,dx
=
\frac{a}{a^2+b^2},
\\
\int_0^\infty
e^{-ax}\sin(bx)\,dx
=
\frac{b}{a^2+b^2}.
\end{gather*}

\begin{gather*}
\int_0^{\pi/2}\sin^2x\,dx
=
\int_0^{\pi/2}\cos^2x\,dx
=
\frac{\pi}{4},
\\
\int_{-\pi/2}^{\pi/2}\cos x\,dx=2.
\end{gather*}

\textbf{Leibniz rule:}
\[
\frac{d}{dt}
\int_{a(t)}^{b(t)}f(x,t)\,dx
=
f(b(t),t)b'(t)
-
f(a(t),t)a'(t)
+
\int_{a(t)}^{b(t)}
\frac{\partial f}{\partial t}\,dx.
\]
\vspace{-12pt}
}{0.46}
\end{mytikzbox}

%=========================================================
% BASIC DIFFERENTIAL EQUATIONS
%=========================================================

\begin{mytikzbox}{Basic differential equations}{
\vspace{-10pt}
\begin{gather*}
Y''+\lambda^2Y=0
\quad\Longrightarrow\quad
Y=A\sin(\lambda x)+B\cos(\lambda x),
\\
Y''-\lambda^2Y=0
\quad\Longrightarrow\quad
Y=A\sinh(\lambda x)+B\cosh(\lambda x)
\\
\hspace{42mm}
=Ce^{\lambda x}+De^{-\lambda x},
\\
Y'+\lambda Y=0
\quad\Longrightarrow\quad
Y=Ae^{-\lambda x},
\\
Y''+\lambda=0
\quad\Longrightarrow\quad
Y=-\frac{\lambda}{2}x^2+Ax+B,
\\
r^2R''+rR'-\lambda^2R=0
\quad\Longrightarrow\quad
R=Ar^\lambda+Br^{-\lambda}.
\end{gather*}

For
\[
y'+P(x)y=Q(x),
\]
\[
y
=
e^{-\int P(x)\,dx}
\left[
\int
e^{\int P(x)\,dx}
Q(x)\,dx+C
\right].
\]
\vspace{-13pt}
}{0.46}
\end{mytikzbox}

%=========================================================
% LAPLACE AND Z TRANSFORMS
%=========================================================

\begin{mytikzbox}{Laplace and $z$ transforms}{
\textbf{Laplace transform:}
\[
\mathcal L\{f(t)\}(s)
=
F(s)
=
\int_0^\infty e^{-st}f(t)\,dt.
\]

\begin{gather*}
\mathcal L\{f'(t)\}
=
sF(s)-f(0),
\\
\mathcal L\{f''(t)\}
=
s^2F(s)-sf(0)-f'(0),
\\
\mathcal L\{e^{at}\}
=
\frac{1}{s-a},
\\
\mathcal L\{\sin(\omega t)\}
=
\frac{\omega}{s^2+\omega^2},
\qquad
\mathcal L\{\cos(\omega t)\}
=
\frac{s}{s^2+\omega^2}.
\end{gather*}

\textbf{$z$ transform:}
\[
\mathcal Z\{f_n\}(z)
=
F(z)
=
\sum_{n=0}^{\infty}f_nz^{-n}.
\]

\[
\mathcal Z\{f_{n+1}\}
=
z\left[F(z)-f_0\right].
\]

For a geometric sequence,
\[
f_n=a^n
\quad\Longrightarrow\quad
F(z)=\frac{z}{z-a}.
\]
}{0.46}
\end{mytikzbox}

%=========================================================
% HUBBARD--STRATONOVICH
%=========================================================

\begin{mytikzbox}{Hubbard--Stratonovich transformation}{
The Gaussian identity
\[
e^{\frac{A}{2}x^2}
=
\frac{1}{\sqrt{2\pi A^{-1}}}
\int_{-\infty}^{\infty}
d\phi\,
e^{-\frac{\phi^2}{2A}+\phi x},
\qquad A>0,
\]
replaces a quadratic interaction by a linear coupling to an
auxiliary field.

Equivalently,
\[
e^{\frac{1}{2}J^TAJ}
=
\frac{1}{\sqrt{\det(2\pi A)}}
\int d^n\phi\,
e^{-\frac{1}{2}\phi^TA^{-1}\phi+J^T\phi}.
\]

Functional form:
\[
e^{\frac{1}{2}\int dx\,dy\,
J(x)A(x,y)J(y)}
\propto
\int\mathcal D\phi\,
e^{
-\frac{1}{2}\int dx\,dy\,
\phi(x)A^{-1}(x,y)\phi(y)
+\int dx\,J(x)\phi(x)
}.
\]

For the Curie--Weiss interaction,
\[
e^{\frac{\beta J}{2N}
\left(\sum_i s_i\right)^2}
\propto
\int dm\,
e^{-\frac{\beta JN}{2}m^2
+\beta Jm\sum_i s_i}.
\]
}{0.46}
\end{mytikzbox}

%=========================================================
% GINZBURG--LANDAU
%=========================================================

\begin{mytikzbox}{Ginzburg--Landau theory}{
For a scalar order parameter $\phi(\mathbf x)$,
\[
\mathcal F[\phi]
=
\int d^dx
\left[
\frac{\kappa}{2}(\nabla\phi)^2
+
\frac{r}{2}\phi^2
+
\frac{u}{4}\phi^4
-
h\phi
\right].
\]

The statistical path integral is
\[
Z[h]
=
\int\mathcal D\phi\,
e^{-\beta\mathcal F[\phi]}.
\]

Mean-field configurations satisfy
\[
\frac{\delta\mathcal F}{\delta\phi}
=
-\kappa\nabla^2\phi
+r\phi
+u\phi^3-h
=
0.
\]

For a uniform field and $h=0$,
\[
r\phi+u\phi^3=0,
\]
and therefore
\[
\phi_0=
\begin{cases}
0, & r>0,\\[2mm]
\displaystyle\pm\sqrt{-\frac{r}{u}},
& r<0.
\end{cases}
\]

Usually,
\[
r=a(T-T_c).
\]
}{0.46}
\end{mytikzbox}

%=========================================================
% GAUSSIAN FIELD THEORY
%=========================================================

\begin{mytikzbox}{Gaussian field theory}{
The quadratic Ginzburg--Landau functional is
\[
\beta\mathcal F_0[\phi]
=
\frac{1}{2}
\int\frac{d^dk}{(2\pi)^d}
\left(r+\kappa k^2\right)
|\phi_{\mathbf k}|^2.
\]

The Fourier convention is
\[
\phi(\mathbf x)
=
\int\frac{d^dk}{(2\pi)^d}
e^{i\mathbf k\cdot\mathbf x}
\phi_{\mathbf k}.
\]

The two-point correlation function is
\[
G(\mathbf k)
=
\langle
\phi_{\mathbf k}\phi_{-\mathbf k}
\rangle
=
\frac{1}{r+\kappa k^2}.
\]

Equivalently,
\[
G(\mathbf x-\mathbf y)
=
\int\frac{d^dk}{(2\pi)^d}
\frac{
e^{i\mathbf k\cdot(\mathbf x-\mathbf y)}
}{
r+\kappa k^2
}.
\]

It satisfies
\[
\left(
-\kappa\nabla^2+r
\right)
G(\mathbf x)
=
\delta^{(d)}(\mathbf x).
\]

The correlation length is
\[
\xi=\sqrt{\frac{\kappa}{r}}.
\]

At large distance,
\[
G(r)
\sim
\frac{e^{-r/\xi}}{r^{(d-1)/2}}.
\]

At criticality, $r=0$,
\[
G(r)\sim\frac{1}{r^{d-2}}.
\]
}{0.46}
\end{mytikzbox}

%=========================================================
% SUSCEPTIBILITY
%=========================================================

\begin{mytikzbox}{Green's functions and susceptibility}{
With a source $h(\mathbf x)$,
\[
Z[h]
=
\int\mathcal D\phi\,
e^{-\beta\mathcal F[\phi]
+\beta\int d^dx\,h(\mathbf x)\phi(\mathbf x)}.
\]

The expectation value is
\[
\langle\phi(\mathbf x)\rangle
=
\frac{1}{\beta}
\frac{\delta\ln Z[h]}
{\delta h(\mathbf x)}.
\]

The connected two-point function is
\[
G_c(\mathbf x,\mathbf y)
=
\langle\phi(\mathbf x)\phi(\mathbf y)\rangle
-
\langle\phi(\mathbf x)\rangle
\langle\phi(\mathbf y)\rangle.
\]

The susceptibility is
\[
\chi(\mathbf x,\mathbf y)
=
\frac{\delta\langle\phi(\mathbf x)\rangle}
{\delta h(\mathbf y)}
=
\beta G_c(\mathbf x,\mathbf y).
\]

For a translationally invariant system,
\[
\chi
=
\int d^dx\,
\chi(\mathbf x)
=
\beta\int d^dx\,G_c(\mathbf x).
\]

In Fourier space,
\[
\chi(\mathbf k)=\beta G_c(\mathbf k).
\]

In the Gaussian approximation,
\[
\chi(\mathbf k)
=
\frac{\beta}{r+\kappa k^2},
\qquad
\chi(0)=\frac{\beta}{r}.
\]
}{0.46}
\end{mytikzbox}

%=========================================================
% O(N) AND GOLDSTONE
%=========================================================

\begin{mytikzbox}{$O(N)$ model and Goldstone modes}{
For an $N$-component order parameter
$\boldsymbol\phi$,
\[
\mathcal F[\boldsymbol\phi]
=
\int d^dx
\left[
\frac{\kappa}{2}
(\nabla\boldsymbol\phi)^2
+
\frac{r}{2}\boldsymbol\phi^2
+
\frac{u}{4}
(\boldsymbol\phi^2)^2
\right].
\]

For $r<0$, choose the vacuum
\[
\boldsymbol\phi_0
=
v\hat{\mathbf e}_1,
\qquad
v^2=-\frac{r}{u}.
\]

Decompose fluctuations as
\[
\boldsymbol\phi
=
(v+\sigma,\pi_1,\ldots,\pi_{N-1}).
\]

To quadratic order,
\[
\mathcal F^{(2)}
=
\frac{1}{2}
\int d^dx
\left[
\kappa(\nabla\sigma)^2
+
m_\sigma^2\sigma^2
+
\kappa
\sum_{a=1}^{N-1}
(\nabla\pi_a)^2
\right],
\]
where
\[
m_\sigma^2=-2r>0,
\qquad
m_{\pi_a}^2=0.
\]

Thus,
\[
G_\sigma(k)
=
\frac{1}{\kappa k^2+m_\sigma^2},
\qquad
G_{\pi}(k)
=
\frac{1}{\kappa k^2}.
\]

There is one Goldstone mode for every broken continuous
symmetry generator.
}{0.46}
\end{mytikzbox}

%=========================================================
% CRITICAL EXPONENTS
%=========================================================

\begin{mytikzbox}{Critical exponents}{
Define the reduced temperature
\[
t=\frac{T-T_c}{T_c}.
\]

Near a continuous phase transition,
\begin{gather*}
C\sim |t|^{-\alpha},
\\
m\sim(-t)^\beta,
\qquad t<0,
\\
\chi\sim |t|^{-\gamma},
\\
m(t=0,h)\sim h^{1/\delta},
\\
\xi\sim |t|^{-\nu},
\\
G(r)\big|_{t=0}
\sim
\frac{1}{r^{d-2+\eta}}.
\end{gather*}

Common scaling relations:
\begin{gather*}
\alpha+2\beta+\gamma=2,
\\
\gamma=\beta(\delta-1),
\\
\gamma=\nu(2-\eta),
\\
2-\alpha=d\nu
\qquad
\text{(hyperscaling)}.
\end{gather*}

Mean-field values:
\[
\alpha=0,
\quad
\beta=\frac{1}{2},
\quad
\gamma=1,
\quad
\delta=3,
\quad
\nu=\frac{1}{2},
\quad
\eta=0.
\]
}{0.46}
\end{mytikzbox}

%=========================================================
% CRITICAL DIMENSIONS
%=========================================================

\begin{mytikzbox}{Upper and lower critical dimensions}{
For the $\phi^4$ or $O(N)$ model, the quartic coupling has
scaling dimension
\[
[u]=4-d.
\]

Therefore,
\[
d_c^{\mathrm{upper}}=4.
\]

For $d>4$, the quartic interaction is irrelevant and
mean-field exponents are valid.

For $d<4$, fluctuations modify the critical exponents.

For continuous $O(N)$ symmetry with $N\geq2$,
Goldstone fluctuations behave as
\[
\langle\pi^2\rangle
\sim
\int\frac{d^dk}{k^2}.
\]

This integral diverges in the infrared for
\[
d\leq2.
\]

Thus,
\[
d_c^{\mathrm{lower}}=2
\qquad
\text{for continuous symmetries}.
\]

For the discrete Ising symmetry,
\[
d_c^{\mathrm{lower}}=1.
\]

The Mermin--Wagner theorem forbids spontaneous breaking
of continuous symmetries at finite temperature for
$d\leq2$ with sufficiently short-range interactions.
}{0.46}
\end{mytikzbox}

%=========================================================
% TRANSFER MATRIX
%=========================================================

\begin{mytikzbox}{Transfer-matrix method}{
For a one-dimensional Ising chain,
\[
H
=
-J\sum_{i=1}^{N}s_is_{i+1}
-h\sum_{i=1}^{N}s_i,
\qquad
s_i=\pm1,
\qquad
s_{N+1}=s_1.
\]

Define the transfer matrix
\[
T_{s,s'}
=
\exp
\left[
\beta Jss'
+
\frac{\beta h}{2}(s+s')
\right].
\]

Explicitly,
\[
T
=
\begin{pmatrix}
e^{\beta J+\beta h} & e^{-\beta J}\\
e^{-\beta J} & e^{\beta J-\beta h}
\end{pmatrix}.
\]

The partition function is
\[
Z_N=\operatorname{Tr}(T^N)
=
\lambda_+^N+\lambda_-^N.
\]

The eigenvalues are
\[
\lambda_\pm
=
e^{\beta J}\cosh(\beta h)
\pm
\sqrt{
e^{2\beta J}\sinh^2(\beta h)
+
e^{-2\beta J}
}.
\]

In the thermodynamic limit,
\[
f
=
-\frac{1}{\beta}
\ln\lambda_+.
\]

For $h=0$,
\[
\lambda_\pm
=
e^{\beta J}\pm e^{-\beta J}.
\]

The spin correlation function is
\[
\langle s_i s_{i+r}\rangle
=
\left(
\frac{\lambda_-}{\lambda_+}
\right)^r
=
\left[
\tanh(\beta J)
\right]^r.
\]

Hence,
\[
\xi^{-1}
=
-\ln[\tanh(\beta J)].
\]
}{0.46}
\end{mytikzbox}

%=========================================================
% CHAPMAN--KOLMOGOROV
%=========================================================

\begin{mytikzbox}{Chapman--Kolmogorov equation}{
For a Markov process,
\[
P(x_3,t_3|x_1,t_1)
=
\int dx_2\,
P(x_3,t_3|x_2,t_2)
P(x_2,t_2|x_1,t_1),
\]
for
\[
t_1<t_2<t_3.
\]

For a time-homogeneous process,
\[
P(x,t+\tau|x_0,t_0)
=
\int dy\,
P(x,\tau|y,0)
P(y,t|x_0,t_0).
\]

The transition probabilities therefore compose according to
\[
\hat P(t+s)=\hat P(t)\hat P(s).
\]
}{0.46}
\end{mytikzbox}

%=========================================================
% KRAMERS--MOYAL AND FOKKER--PLANCK
%=========================================================

\begin{mytikzbox}{Kramers--Moyal expansion}{
Define the conditional jump moments
\[
M^{(n)}(x,t;\tau)
=
\frac{1}{\tau}
\int d(\Delta x)\,
(\Delta x)^n
P(x+\Delta x,t+\tau|x,t).
\]

The Kramers--Moyal coefficients are
\[
D^{(n)}(x,t)
=
\frac{1}{n!}
\lim_{\tau\to0}
M^{(n)}(x,t;\tau).
\]

The probability density obeys
\[
\frac{\partial P}{\partial t}
=
\sum_{n=1}^{\infty}
(-\partial_x)^n
\left[
D^{(n)}(x,t)P(x,t)
\right].
\]

If
\[
D^{(n)}=0
\qquad
\text{for all }n\geq3,
\]
the expansion truncates to the Fokker--Planck equation:
\[
\frac{\partial P}{\partial t}
=
-\partial_x
\left[
D^{(1)}P
\right]
+
\partial_x^2
\left[
D^{(2)}P
\right].
\]

Pawula's theorem: a finite truncation beyond second order is
generally inconsistent with positivity.
}{0.46}
\end{mytikzbox}

%=========================================================
% LANGEVIN TO FOKKER--PLANCK
%=========================================================

\begin{mytikzbox}{Langevin and Fokker--Planck equations}{
For the Itô stochastic differential equation
\[
dx
=
A(x,t)\,dt
+
B(x,t)\,dW_t,
\]
with
\[
\langle dW_t\rangle=0,
\qquad
\langle dW_t^2\rangle=dt,
\]
the Fokker--Planck equation is
\[
\frac{\partial P}{\partial t}
=
-\frac{\partial}{\partial x}
\left[
A(x,t)P
\right]
+
\frac{1}{2}
\frac{\partial^2}{\partial x^2}
\left[
B^2(x,t)P
\right].
\]

Thus,
\[
D^{(1)}=A,
\qquad
D^{(2)}=\frac{B^2}{2}.
\]

For additive noise,
\[
\dot x=A(x)+\eta(t),
\]
\[
\langle\eta(t)\rangle=0,
\qquad
\langle\eta(t)\eta(t')\rangle
=
2D\delta(t-t'),
\]
one obtains
\[
\frac{\partial P}{\partial t}
=
-\partial_x(AP)+D\partial_x^2P.
\]
}{0.46}
\end{mytikzbox}

%=========================================================
% MASTER EQUATION
%=========================================================

\begin{mytikzbox}{Master equation}{
For discrete states $n$,
\[
\frac{dP_n}{dt}
=
\sum_m
\left[
W_{n\leftarrow m}P_m
-
W_{m\leftarrow n}P_n
\right].
\]

Probability conservation:
\[
\sum_nP_n=1,
\qquad
\sum_n\frac{dP_n}{dt}=0.
\]

A stationary distribution satisfies
\[
\frac{dP_n^{\mathrm{st}}}{dt}=0.
\]

Detailed balance is the stronger condition
\[
W_{n\leftarrow m}P_m^{\mathrm{eq}}
=
W_{m\leftarrow n}P_n^{\mathrm{eq}}
\]
for every pair $n,m$.

Detailed balance implies stationarity, but stationarity does
not necessarily imply detailed balance.
}{0.46}
\end{mytikzbox}

%=========================================================
% METROPOLIS
%=========================================================

\begin{mytikzbox}{Metropolis algorithm}{
For a configuration $C$ with energy $E(C)$, propose
\[
C\longrightarrow C'.
\]

The energy difference is
\[
\Delta E=E(C')-E(C).
\]

For a symmetric proposal probability, accept with probability
\[
A(C\to C')
=
\min\left\{
1,e^{-\beta\Delta E}
\right\}.
\]

Equivalently,
\[
\Delta E\leq0
\quad\Longrightarrow\quad
\text{always accept},
\]
\[
\Delta E>0
\quad\Longrightarrow\quad
\text{accept with probability }
e^{-\beta\Delta E}.
\]

This choice satisfies detailed balance with
\[
P_{\mathrm{eq}}(C)
=
\frac{e^{-\beta E(C)}}{Z}.
\]
}{0.46}
\end{mytikzbox}

%=========================================================
% CLUSTER AND VIRIAL EXPANSIONS
%=========================================================

\begin{mytikzbox}{Cluster and virial expansions}{
For a classical interacting gas,
\[
e^{-\beta U}
=
\prod_{i<j}
\left(1+f_{ij}\right),
\]
where the Mayer function is
\[
f_{ij}
=
e^{-\beta u(r_{ij})}-1.
\]

The pressure has the activity expansion
\[
\beta P
=
\sum_{\ell=1}^{\infty}
b_\ell z^\ell,
\]
while the density is
\[
\rho
=
\sum_{\ell=1}^{\infty}
\ell b_\ell z^\ell.
\]

The virial expansion is
\[
\frac{\beta P}{\rho}
=
1+B_2\rho+B_3\rho^2+\cdots.
\]

Relations between cluster and virial coefficients:
\[
B_2=-\frac{b_2}{b_1^2},
\]
\[
B_3
=
\frac{4b_2^2}{b_1^4}
-
\frac{2b_3}{b_1^3}.
\]

With the common convention $b_1=1$,
\[
B_2=-b_2,
\qquad
B_3=4b_2^2-2b_3.
\]

For a pair potential,
\[
B_2
=
-\frac{1}{2}
\int d^dr\,
\left[
e^{-\beta u(r)}-1
\right].
\]
}{0.46}
\end{mytikzbox}
